\section{The wing as a stated choice}
\label{sec:vswingchoice}

Inside the envelope the data has no further opinion, yet every fitted
model draws exactly one curve there.  What selects the curve is the
model's structure --- decisions made when the family was designed,
long before this node's quotes arrived.  This section reads those
decisions off the three families, states them as contracts, and then
demolishes a comfortable inference about them.

Recall the three families' wings; each was derived in its own chapter,
and here they only need to be set side by side.  The LQD slice of
Chapter~2 has exponential log-return tails \emph{by construction}, and
its wing slopes are closed functions of its two endpoint speeds
(\cref{prop:leeclosed}): admissibility is structural, no cap or check
required.  The SVI slice of Chapter~3 has exactly straight
total-variance wings, $\beta_{L,R}=b_S(1\mp\rho_S)$ --- the
$S$-subscripts mark Chapter~3's raw SVI parameters
(\cref{prop:svigeom}) --- and its calibration fences them with a cap
set strictly inside the ceiling (\cref{eq:leecap}): strictly, because
\cref{sec:vsceiling} showed the boundary ray itself is inadmissible.  The MCS
slice of Chapter~3 splits the job: its base owns the wings
(\cref{eq:mcsbetak}), and its hats are built to vanish at large
strikes, so adding hats cannot move the slopes --- but nothing in the
family \emph{forces} the base's slopes into the admissible cone;
their admissibility is diagnosed after the fit, not enforced during
it.  On the running node the three contracts print the slopes of
\cref{tab:vswings}: every wing far below the ceiling, and the put
slopes several times the call slopes.  That asymmetry is the equity
skew --- the standing pattern that downside strikes carry more
implied variance than upside ones --- read here in its asymptotic
form.

\begin{table}[htbp]
\centering
\caption{The three wing contracts on the running node.  Slopes are
asymptotic total-variance wing slopes; all sit far inside the ceiling
of \cref{prop:vsceiling}.}
\label{tab:vswings}
\begin{tabular}{lccl}
\toprule
Family & $\beta_L$ (put) & $\beta_R$ (call) & how the wing is governed\\
\midrule
LQD & \MacWingBetaLqdL & \MacWingBetaLqdR &
structural: exponential tails, slopes from \cref{prop:leeclosed}\\
SVI & \MacWingBetaSviL & \MacWingBetaSviR &
capped: straight wings under \cref{eq:leecap}\\
MCS & \MacWingBetaMcsL & \MacWingBetaMcsR &
inherited: base wings, hats neutral; diagnosed, not enforced\\
\bottomrule
\end{tabular}
\end{table}

Three families, three slopes per side, all admissible --- and all
different, as \cref{fig:vsenvelope}(b) drew.  It bears saying plainly:
past the last quote, the divergence between the families is not an
error to be averaged away.  Choosing a family \emph{is} choosing an
extrapolation law.  What a careful fitter owes its user is not
agreement in the wings --- nothing can supply that --- but a stated
contract: which curve will be drawn, governed by what mechanism, with
what guarantee.

There is, however, a gap in those guarantees, and it is the section's
second lesson.  All three contracts of \cref{tab:vswings} speak about
\emph{slopes} --- asymptotic statements, properties of the limit
$k\to\infty$.  The comfortable inference is: admissible slopes,
therefore admissible wing.  It is false, and one deliberate
construction shows why.

\begin{figure}[htbp]
\centering
\paperfig{\textwidth}{fig_wing_limit.pdf}
\caption{The limit is not the wing.  One zero-wing hat is added to
the fitted MCS slice, centered beyond the last put-side quote.
(a)~Total variance: the two curves are indistinguishable except for
the small bump at the hat; the measured far-field wing slopes change
by \MacWingHatSlopeDiff\ --- an exact zero in double precision.
(b)~The density factor $g_{\rm D}$: the hat drives it to
\MacWingHatGmin\ at $k=\MacWingHatKmin$, negative density at a finite
strike, while the fitted slice (solid) stays nonnegative to rounding
(worst \MacWingHatBaseGmin).}
\label{fig:winglimit}
\end{figure}

Take the fitted MCS slice of the running node and add one more of its
own hat kernels (\cref{eq:hatdef}) --- amplitude negative, centered
just beyond the last put-side quote, where quotes no longer object.
\Cref{fig:winglimit}(a) shows the smile before and after: the two
curves coincide to the eye everywhere except a small dent near
$k=\MacWingHatKmin$, and the wings are untouched --- measuring the
realized slope far out, the change is \MacWingHatSlopeDiff, an exact
zero at double precision.  The inheritance contract has done its job
perfectly.  Panel~(b) shows what the contract does not cover.  Recall
the Durrleman factor $g_{\rm D}$ from \cref{eq:durrleman}: up to a
positive factor it is the probability density the smile implies, so a
valid smile needs $g_{\rm D}\ge0$ everywhere.  The fitted slice
(solid) keeps it nonnegative across the range --- its worst value is
\MacWingHatBaseGmin, zero at rounding scale.  The hatted slice
(dashed) plunges to $g_{\rm D}=\MacWingHatGmin$ at
$k=\MacWingHatKmin$: negative probability, a butterfly arbitrage,
parked at a finite strike in the extrapolated region --- with both
asymptotic slopes exactly as admissible as before.  The dent sits in
the quiet stretch just outside the quoted span, precisely where no
quote penalizes it and no slope condition sees it.  Asymptotic
admissibility and finite-strike admissibility are independent, and
the second needs its own police.

The demonstration settles how the extrapolated region must be
policed, and the rule is worth stating because it cuts both ways.
Conditions that are properties of a \emph{single} curve --- density
positivity above, call convexity --- must be checked \emph{into} the
wings: one law is being asserted there, the condition needs no second
curve to compare against, and abandoning it past the last quote is
exactly how the hat's dent would go unnoticed.  Conditions that
\emph{compare two} curves are different.  Chapter~2's calendar
control (\cref{eq:calrows}) already made the choice: it enforces the
order between neighboring expiries only on the span where quotes pin
both curves.  Beyond that span, both curves are chosen conventions
--- two entries of \cref{tab:vswings}, possibly from different rows.
A calendar inversion there --- the later expiry's curve dipping below
the earlier one's at some unquoted strike --- is a disagreement
between two fictions, not evidence of mispricing; enforcing the order
there would bend the fits where the data \emph{does} live in order to
reconcile them where it does not.  One
curve: extend the check.  Two curves: confine it to common data.

The chapter's statics are now complete for a single expiry: the
number, its three computations, its unquoted share, the proved
envelope, and the stated choice within it.  One dimension remains
--- maturity.
